\section{Testing both minuscule directions}
\label{sec:kt-dual-hecke}

A trace comparison in rank three does not by itself impose the
inverse-trace condition.  The stored harmonic model permits both
conditions to be tested at the same coefficient prime.  The following
identity specifies that test without identifying a coefficient prime
with its complex conjugate.

\begin{lemma}[KT-U3-DUAL-MATRIX]
Let $C$ interchange the two sets of polynomial variables in the
harmonic model, and choose the saved basis with
$C=\operatorname{diag}(s_i)$, $s_i\in\{1,-1\}$.
Write $M_p=M_{p,0}+\omega M_{p,1}$ for the scaled negative
minuscule operator.  Put $\delta=1+2\omega$, $\delta^2=-3$,
and let $B$ be diagonal with entry $1$ on the plus basis
vectors and $\delta$ on the minus basis vectors.  In the
rational coefficient basis used in the data, the two scaled
minuscule matrices are
\[
       R_p=B^{-1}M_pB,\qquad R_p^{\rm opp}=CR_pC.
                                                               \tag{KT-D.1}
\]
At any modulus for which all recorded matrix denominators and
$6$ are units, the simultaneous kernel at a rational target $t_p$
is the direct sum of its two $C$-parity kernels.  On a vector of
one parity, it is enough to impose $R_p-t_pI$ for each $p$.
\end{lemma}
\begin{proof}
The full harmonic action is
$\rho(g)f(z,w)=f(g^{-1}z,\overline g^{-1}w)$.
Consequently $C\rho(g)C=\rho(\overline g)$.
Conjugation interchanges the positive and negative minuscule
coset sets, so the opposite matrix is $CM_pC$ on this same
space.  Since the saved basis is rational, it also equals
$\overline{M_p}$, giving
\[
 CM_{p,0}C=M_{p,0}-M_{p,1},\qquad
 CM_{p,1}C=-M_{p,1}.
\]
For equal-parity indices the $\omega$ coefficient is zero.
For unequal-parity indices it is twice the rational coefficient,
so that entry of $M_p$ is $\delta(M_{p,0})_{ij}$.
Conjugating by $B$ makes such an entry
$-3(M_{p,0})_{ij}$ in a plus row and
$(M_{p,0})_{ij}$ in a minus row.  This is exactly the saved
rational matrix.  Since $B$ commutes with $C$, it also proves
the second formula in (KT-D.1).

These are equalities before reduction.  At a coefficient
prime away from three, $B$ is invertible; both rational
matrices reduce in one and the same residue field.  No action
of complex conjugation on that prime is assumed.
The common kernel of $R_p-t_pI$ and $CR_pC-t_pI$, over all
chosen $p$, is invariant under $C$.  Since two is a unit,
the projectors $(1\pm C)/2$ split it into its parity parts.
For $Cv=sv$, the identity
\[
        (CR_pC-t_pI)v=sC(R_p-t_pI)v
\]
shows that the opposite condition is automatic after the
first one.  Restricting the columns of the stacked matrices
$R_p-t_pI$ to each parity therefore computes the exact
simultaneous kernel.
\end{proof}

\paragraph{Exact computation at the three smaller residual moduli.}
The program \texttt{check\_dual\_hecke.py} verifies all coefficient
identities above against the complete stored harmonic matrices,
and checks every denominator for invertibility modulo the
specified modulus.  Its elimination divides only by units.
The target at $p$ is $t_p=1+p^e+p^{2e}$.
For $p=7$ alone, for $p=13$ alone, and for the two primes
together, the parity ranks are the same:
\[
\begin{array}{c|r|c|c}
e&\text{modulus}&\text{plus rank/dimension}&
                         \text{minus rank/dimension}\\\hline
26&697446897311&44/44&39/40\\
28&323026426255399&55/55&49/50\\
30&236017358500033404991&65/65&59/60
\end{array}
\]
Every reduced matrix has zero residual rows after the listed
unit pivots.  In each minus block the program constructs a
kernel vector with a coordinate equal to one and verifies it
directly in all four full matrices.  Thus the vector remains
nonzero at every prime divisor of the listed modulus.
The complete pivot data and primitive kernel vectors are in
\texttt{dual\_hecke\_check.json}.  No primality test is needed
for these assertions.

The result is a failure of this finite Hecke exclusion test:
adding the opposite minuscule generators does not eliminate
these three residual obstructions.  It does not assert that
every resulting finite-prime congruence is an admissible
deformation-theoretic source.  KT-U3-SYM2-INJECT and the exact
factor checks locate the smaller-modulus spectral support in
the constructed symmetric-square family; both directions in
that family necessarily have equal eigenvalues.
